\section{Companion lab: extend arithmetic beyond one register}
\label{revision-lab:rsa}

Word operations have a fixed width. When a result needs more bits, an assembly
program must keep those bits explicitly. This companion develops unsigned
multiword arithmetic through a small RSA calculation. It revises original project
\texttt{027-rsa-arithmetic}; the numerical interfaces also work for composite
moduli unrelated to an RSA key.

Use the native lesson \emph{Extend arithmetic beyond one register}. Implement
the seven entries in its \texttt{independent/rsa-arithmetic.asm}. The same bare
starter is retained at \texttt{revision/labs/files/rsa-arithmetic.asm} for
standalone work. The skeleton supplies caller
objects and an aligned stack area, with a deliberate \texttt{break} at each
unfinished function. The original package retains its cumulative guided source
and public checks. Your completed functions can choose different helpers,
registers, frames and exponent scans. Start with the scalar cases before extending
the representation to four limbs.

The native lesson begins with short worked sources in its \texttt{worked}
directory. Each exposes one operation before asking you to implement a complete
interface. Answer the inline prediction, inspect the requested Debug boundaries,
then return to the continuing independent source. Finish the Debug process before
another launch or source edit. The worked files remain unchanged throughout the
required observations; there is no need to change your implementation's mode
word to follow them.

\subsection*{Use the gcd to explain a small exponent pair}

For positive modulus $N$, the \emph{residue} of $x$ is the remainder in the
range zero through $N-1$. Two integers are congruent modulo $N$ when their
difference is divisible by $N$. Reducing intermediate values to residues keeps
the numbers smaller without changing the final modular result.

Two integers are coprime when their greatest common divisor is one.
For an RSA arithmetic example, choose primes 61 and 53. Their product is
$N=3233$. The count of residues coprime to this product is
$(61-1)(53-1)=3120$. Choose public exponent 17. Its greatest common divisor
with 3120 is one, which makes an inverse modulo 3120 possible.

\begin{table}[!ht]
  \centering
  \begin{tabular}{r r r r}
    \toprule
    $A$  & $B$ & Quotient & Remainder \\
    \midrule
    3120 & 17  & 183      & 9         \\
    17   & 9   & 1        & 8         \\
    9    & 8   & 1        & 1         \\
    8    & 1   & 8        & 0         \\
    \bottomrule
  \end{tabular}
  \caption{Each Euclidean iteration replaces $(A,B)$ with $(B,A\bmod B)$. The last nonzero remainder is one. The terminating zero is not the gcd.}
  \label{tab:revision-rsa-euclid}
\end{table}


\texttt{rsa\_gcd64(a0 A,a1 B)} accepts full unsigned 64-bit values.
At each iteration, compute the unsigned remainder and replace both operands.
The loop stops when $B$ is zero, returning the current $A$. Thus
$\gcd(A,0)=A$; our local convention also defines $\gcd(0,0)=0$.
Check for a zero divisor before computing the remainder.

Begin the native observation with \texttt{worked/remainders.asm}. Its first
pair is 13,11, whose unsigned remainder is two. Before replacing either operand,
identify which old value each move still needs. The next pair must be 11,2;
its remainder is one. Figure~\ref{fig:revision-rsa-native-pause} shows that first
pause. The blue execution row and selected frame describe the current machine
state. The yellow lesson pointer identifies the next requested stop after the
lesson has accepted the current observation.

\begin{figure}[!ht]
\centering
\begin{tikzpicture}[alt={Three crops from one actual light-theme IDE photograph of the public remainders experiment. The blue execution row is the move at line 9, confirmed by the selected paused frame. Registers t0 and t1 still contain decimal 13 and 11, while t2 contains their unsigned remainder two. The yellow lesson pointer has moved to line 13, the next requested stop after updating the pair and computing another remainder. Label A identifies the current pause and label B identifies the next lesson target.}]
\node[anchor=north west,inner sep=0,font=\small\bfseries] (title) at (0,0) {Read the remainder before updating the pair};
\node[anchor=north west,inner sep=0] (code) at ([yshift=-2mm]title.south west) {\MARSRevisionCrop{revision/captures/rsa/rsa-native-pause.png}{0.96\linewidth}{0.42968750}{0.53777778}{0.17968750}{0.18888889}};
\begin{scope}[shift={(code.south west)},x={($(code.south east)-(code.south west)$)},y={($(code.north west)-(code.south west)$)}]
\draw[black,line width=.7pt] (.11,.5528) rectangle (.995,.6423);
\draw[black,line width=.7pt] (.11,.1951) rectangle (.995,.2846);
\node[draw=black,fill=white,inner sep=1pt,font=\small\bfseries] at (.075,.60) {A};
\node[draw=black,fill=white,inner sep=1pt,font=\small\bfseries] at (.075,.24) {B};
\end{scope}
\node[anchor=north west,inner sep=0] (frame) at ([yshift=-2mm]code.south west) {\MARSRevisionCrop{revision/captures/rsa/rsa-native-pause.png}{0.4416\linewidth}{0.03750000}{0.22222222}{0.78281250}{0.75111111}};
\node[anchor=north west,inner sep=0] (value) at ([yshift=-2mm]frame.south west) {\MARSRevisionCrop{revision/captures/rsa/rsa-native-pause.png}{0.768\linewidth}{0.24765625}{0.03333333}{0.43984375}{0.89777778}};
\end{tikzpicture}
\caption{At the current pause (A), \texttt{t0}=13 and \texttt{t1}=11 have
produced \texttt{t2}=2. The register view is hexadecimal: \texttt{0d} and
\texttt{0b} mean decimal 13 and 11. The move at line 9 has not executed. The lesson accepts
that observation and points to the next requested boundary at line 13 (B).
Place the next breakpoint there and use Resume: the pair becomes 11,2, and the
next unsigned remainder is one. The frame and register crops come from the same
photograph. The program is still paused at A.}
\label{fig:revision-rsa-native-pause}
\end{figure}


The final part of this worked source divides the all-one doubleword by 17.
Unsigned interpretation gives $2^{64}-1$, with remainder zero; signed
interpretation gives $-1$, with signed remainder $-1$. The register bits alone
do not choose the operation. Use the interpretation promised by the interface,
then complete the gcd assessment before moving on to modular powers.

Work backward through the 3120,17 calculation:
\[
 1=9-8=2\cdot9-17=2\cdot3120-367\cdot17.
\]
Consequently $-367$ is an inverse of 17 modulo 3120. Its nonnegative residue is
2753, and $17\cdot2753=15\cdot3120+1$. This derives the exponent used below;
it is not a second constant to memorize. You need not implement key generation
or an extended-gcd interface for this assignment.

The gcd entry returns its numerical result in \texttt{v0} and success status
zero in \texttt{v1}. Other entries return status in \texttt{v0} and write their
numerical result to memory. The supplied caller pauses at \texttt{gcd64\_return}
before moving the gcd status into \texttt{v0}. Inspect the gcd value at that
first boundary; \texttt{main\_done} is too late to expect it in \texttt{v0}.
A function returning a high-bit-set gcd must preserve all 64 result bits.

\Needspace{11\baselineskip}
\subsection*{Reduce a power without constructing the full integer}

Take message integer 65. The small forward operation is
$65^{17}\bmod3233=2790$, and the reverse calculation is
$2790^{2753}\bmod3233=65$. Euler's theorem gives
$65^{3120}\equiv1\pmod{3233}$ because 65 is coprime to 3233. Since
$17\cdot2753=1+15\cdot3120$, the combined operation is
$65(65^{3120})^{15}$, whose residue is 65.
The program will compute the remainders without ever storing $65^{17}$ itself.

\texttt{rsa\_pow32(a0 base,a1 exponent,a2 modulus,a3 output)} receives
zero-extended unsigned word operands. Its successful output is one little-endian
word. Start with result $r=1\bmod N$, current base $x=base\bmod N$ and remaining
exponent $y=exponent$. At every loop boundary,
\[
 r\,x^y\equiv \mathit{base}^{\mathit{exponent}}\pmod N.
\]
If $y$ is odd, multiply $r$ by $x$ and reduce. Shift $y$ right logically by one,
and square and reduce $x$ if another exponent bit remains. Each iteration
removes one bit, rather than performing one multiplication for each unit of the
original exponent.

\begin{table}[!ht]
  \centering
  \begin{tabular}{r r r r}
    \toprule
    Exponent before & Result after & Base after & Exponent after \\
    \midrule
    17              & 65           & 992        & 8              \\
    8               & 65           & 1232       & 4              \\
    4               & 65           & 1547       & 2              \\
    2               & 65           & 789        & 1              \\
    1               & 2790         & 789        & 0              \\
    \bottomrule
  \end{tabular}
  \caption{Least-significant-bit-first exponentiation of 65 to power 17 modulo 3233. Initially the result is one and base is 65. This version stops before squaring an unused base after the last exponent bit.}
  \label{tab:revision-rsa-power}
\end{table}


Predict the five rows before running. Explain why the result stays 65 while
remaining exponents eight, four and two are processed. Those iterations still
change the current base; skipping their squaring would make the final odd bit
multiply by the wrong power.

\begin{figure}[!ht]
\centering
\begin{tikzpicture}[alt={Actual native Memory row after scalar modular exponentiation. The first four bytes are e6 0a 00 00, encoding 2790, followed by twelve unchanged a5 bytes. An outline identifies the four-byte output.}]
\node[anchor=north west,inner sep=0,font=\small\bfseries] (title) at (0,0) {One word holds the reduced power};
\node[anchor=north west,inner sep=0] (a) at ([yshift=-2mm]title.south west) {\MARSRevisionCrop{revision/captures/rsa/rsa-word.png}{0.99\linewidth}{0.03125}{0.33888889}{0.58203125}{0.60888889}};
\begin{scope}[shift={(a.south west)},x={($(a.south east)-(a.south west)$)},y={($(a.north west)-(a.south west)$)}]
\draw[black,line width=.7pt] (.20,.01) rectangle (.39,.467);
\end{scope}
\end{tikzpicture}
\caption{The outlined bytes e6 0a 00 00 encode decimal 2790, the residue of 65 to power 17 modulo 3233. The following twelve pictured bytes remain a5 and are outside this four-byte result.}
\label{fig:revision-rsa-native-word}
\end{figure}


\Needspace{11\baselineskip}
Two word residues can have a product requiring 64 bits. Use a doubleword
multiply and an unsigned doubleword remainder for this path:

\MARSCodeLabel{Keep the entire product until after reduction}
\begin{minted}{mars-idea}
    dmulu $t1, $t1, $t0
    dmodu $t1, $t1, $a2
\end{minted}

Here \texttt{t1} is the accumulated residue, \texttt{t0} the current base and
\texttt{a2} the zero-extended modulus. A word multiply discards needed high bits.
A signed doubleword remainder also fails when the complete product's top bit is
one. For base \texttt{fffffffe}, exponent two and modulus \texttt{ffffffff},
the product is \texttt{fffffffc00000004} and the correct residue is one.
Interpreting those product bits as a negative signed integer gives a different
remainder. Use \texttt{lwu} when the caller loads unsigned word operands.

Exponent zero returns $1\bmod N$, including for a zero base. For modulus one,
that answer is zero. Modulus zero is an error: return status one without
accessing the output. Test these cases separately from a large positive exponent.

In \texttt{worked/word-product.asm}, compare \texttt{lw} and \texttt{lwu}
on the same all-one limb. Then inspect the complete product
\texttt{fffffffe00000001}. Unsigned reduction modulo \texttt{ffffffff}
gives zero, while its low word is one and signed reduction gives $-1$. The
experiment separates three values that can look plausible when printed in
hexadecimal. Finish the scalar-power checks with the complete unsigned dividend.

\subsection*{Represent a wider number with unsigned limbs}

Let $\beta=2^{32}$. Four \emph{limbs}, each an unsigned word, represent
\[
 A=A_0+A_1\beta+A_2\beta^2+A_3\beta^3.
\]
The least significant limb is at the lowest address, and the bytes inside each
limb are little-endian. This is the local arithmetic representation; it differs
from the big-endian byte strings used by the RSA conversion primitives in
\href{https://www.rfc-editor.org/rfc/rfc8017.html}{RFC 8017}.
Do not treat a dump of these limbs as an encoded RSA message.

Unsigned loads matter even when every stored limb is a perfectly valid word.
Loading \texttt{ffffffff} with \texttt{lwu} produces the value $\beta-1$ in a
64-bit register. Sign extension would instead change the upper bits used by
later doubleword multiplication. Pointer arithmetic remains full width regardless
of the limb width.

\Needspace{7\baselineskip}
\texttt{rsa\_mul128(a0 A,a1 B,a2 output)} writes the full eight-limb product.
One schoolbook implementation clears all eight output words, then processes one
input limb against each limb of the other input. For an inner product, add
$A_iB_j$, the existing output limb and the incoming carry. Store the low word
and retain the high word as the next carry.

This ordering has a useful bound. Each operand, existing limb and carry is at
most $\beta-1$, so the inner total is at most
\[
 (\beta-1)^2+(\beta-1)+(\beta-1)=\beta^2-1.
\]
It fits exactly within an unsigned doubleword. Summing several unreduced limb
products at once would require a different bound and possibly more temporary
bits. The reason this particular loop fits is its carry discipline.

For a hand-checkable example, $(2\beta-1)^2$ has low-to-high limbs
$1,\beta-4,3,0,0,0,0,0$. For the largest input $M=\beta^4-1$, its square has limbs
$1,0,0,0,\beta-2,\beta-1,\beta-1,\beta-1$. Check all eight words and the bytes immediately
after the selected 32-byte result. Returning only the low four words implements
a different operation.

\begin{figure}[!ht]
  \centering
  \begin{tikzpicture}[alt={Actual native Memory view of the complete thirty-two-byte product of two maximum unsigned 128-bit values. First-row words are one and three zeros. Second-row words are fffffffe followed by three ffffffff words. An outline encloses all eight limbs.}]
    \node[anchor=north west,inner sep=0,font=\small\bfseries] (title) at (0,0) {Keep all eight product limbs};
    \node[anchor=north west,inner sep=0] (a) at ([yshift=-2mm]title.south west) {\MARSRevisionCrop{revision/captures/rsa/rsa-product.png}{0.99\linewidth}{0.03125}{0.31555556}{0.58203125}{0.60888889}};
    \begin{scope}[shift={(a.south west)},x={($(a.south east)-(a.south west)$)},y={($(a.north west)-(a.south west)$)}]
      \draw[black,line width=.7pt] (.20,.005) rectangle (.965,.6176);
    \end{scope}
  \end{tikzpicture}
  \caption{The full product of two maximum unsigned 128-bit inputs occupies the outlined 32 bytes. Low-to-high limbs are one, zero, zero, zero, fffffffe, ffffffff, ffffffff and ffffffff.}
  \label{fig:revision-rsa-native-product}
\end{figure}


The native \texttt{worked/limb-column.asm} adds two maximum limbs and a
supplied carry of one. Its 33-bit sum is \texttt{1ffffffff}: low limb
\texttt{ffffffff}, outgoing carry one. This isolates one column. In a full
multiplication, several limb products can contribute to the same position;
account for those contributions before declaring that output limb finished.
The independent full-product checks also require zero high limbs and squaring
through one shared read-only input object.

\subsection*{Keep the carry that lies beyond the output array}

\texttt{rsa\_addmod128(a0 A,a1 B,a2 N,a3 output)} requires reduced inputs
$A<N$ and $B<N$. Their sum is below $2N$, so at most one subtraction of $N$
is needed. However, the unreduced sum can require 129 bits. Four stored result
limbs alone cannot say whether that subtraction is necessary.

Add from low limbs toward high limbs, retaining a carry into each next position.
After the fourth limb, retain the final carry separately. If it is one, subtract
$N$ even when the stored low 128 bits compare below $N$. Otherwise compare the
four-limb sum with $N$, starting at the most significant limb, and subtract when
the sum is at least $N$. Equality must produce zero.

Consider $N=2^{128}-159$ and $A=B=N-1$. The mathematical sum is
$2^{129}-320$. Its low 128 bits represent $2^{128}-320$, which is below $N$,
but its carry is one. The correct reduced result is $N-2=2^{128}-161$.
A comparison that ignores the carry would return the wrong provisional sum.
When subtracting $N$ from the low limbs, the borrow consumes the retained high
carry; the final reduced value fits in four limbs.

\begin{figure}[!ht]
  \centering
  \begin{tikzpicture}[alt={Actual native Memory row during modular addition, before conditional subtraction. The low limbs begin c0 fe ff ff and otherwise contain ff. An outline encloses these provisional sixteen bytes. The separately verified 129th carry bit is not pictured.}]
    \node[anchor=north west,inner sep=0,font=\small\bfseries] (title) at (0,0) {The provisional sum still needs subtraction};
    \node[anchor=north west,inner sep=0] (a) at ([yshift=-2mm]title.south west) {\MARSRevisionCrop{revision/captures/rsa/rsa-carry.png}{0.99\linewidth}{0.03125}{0.33888889}{0.58203125}{0.60888889}};
    \begin{scope}[shift={(a.south west)},x={($(a.south east)-(a.south west)$)},y={($(a.north west)-(a.south west)$)}]
      \draw[black,line width=.7pt] (.20,.01) rectangle (.965,.467);
    \end{scope}
  \end{tikzpicture}
  \caption{The pictured low 128 bits are the provisional sum, beginning c0 fe ff ff. Its retained high carry is one, verified separately and not pictured. Despite these low limbs being below N, subtraction must produce N minus two before return.}
  \label{fig:revision-rsa-native-carry}
\end{figure}


\Needspace{7\baselineskip}
Subtraction also proceeds from low to high, with an incoming borrow of zero
or one. For zero-extended word operands, a negative limb difference occupies
the high unsigned doubleword range; extract the borrow before retaining only
the low word. Test a borrow that propagates across several zero limbs.

This entry permits output to exactly alias either or both operands. A limbwise
implementation can load both current operand limbs before replacing the current
output limb, provided no later step needs those original limbs. The modulus
stays disjoint from output. An arbitrary partial overlap has no such guarantee;
a store could replace a higher source limb before it is read.

Use \texttt{worked/carry-reduction.asm} to inspect the decision with an
eight-bit model: $N=251$ and $A=B=250$. The full sum is 500, its low eight bits
are 244, and the outgoing carry is one. Comparing only 244 with 251 would skip
a necessary subtraction. The result is $500-251=249$. Transfer that reasoning
to the four-limb representation; the small experiment can display a complete
intermediate in one register, while your general implementation must retain its
extra carry explicitly.

\subsection*{Reduce arbitrary inputs before using residue-only helpers}

\texttt{rsa\_reduce128(a0 A,a1 N,a2 output)} accepts any unsigned 128-bit
input, not merely a value less than twice $N$. Repeatedly subtracting a small
modulus could require an impractical number of iterations. Instead, visit the
input bits from most significant to least significant.

Maintain the residue $r$ of the prefix already visited. For next bit $b$, the
new prefix has residue $(2r+b)\bmod N$. Since $r<N$, this step has a bounded
amount of reduction work. For binary input \texttt{1101}, decimal 13, modulo
five, the successive residues are one, three, one and three.
The order of the bits matters: low-to-high scanning does not follow this prefix
recurrence.

One implementation uses the reduced addition helper to compute $r+r$ and then,
when the incoming bit is one, to add one. Handle modulus one first: a literal
one is not a reduced operand for that modulus. For $N>1$, both zero and one
are valid residues. Initialize the running residue independently of any old
output or private workspace contents.

\Needspace{9\baselineskip}
\texttt{rsa\_mulmod128(a0 A,a1 B,a2 N,a3 output)} accepts arbitrary inputs
too. Reduce its current multiplicand before passing it to addmod. Keep accumulator
$r=0$, reduced current value $x=A\bmod N$ and remaining multiplier $y=B$.
When the low bit of $y$ is one, replace $r$ by $(r+x)\bmod N$. Shift $y$
right and double $x$ modulo $N$ for the next bit. The invariant is
$r+xy\equiv AB\pmod N$.

\begin{table}[!ht]
\centering
\begin{tabular}{r r r}
\toprule
Accumulator after & Current residue after & Multiplier after\\
\midrule
13 & 9 & 5\\
5 & 1 & 2\\
5 & 2 & 1\\
7 & 2 & 0\\
\bottomrule
\end{tabular}
\caption{Modular multiplication of 13 by 11 modulo 17. Initially the accumulator is zero, current residue 13 and multiplier 11. The final answer is seven; no unused last doubling is performed.}
\label{tab:revision-rsa-modproduct}
\end{table}


This computes a modular product without first constructing a 256-bit temporary
product. Both approaches can be correct, but a full-product approach also needs
a reduction algorithm that accepts all 256 product bits. Passing only the low
128 bits to \texttt{rsa\_reduce128} loses information.

Trace \texttt{worked/prefix-reduction.asm} through all four bits of
\texttt{1101} modulo five. Its returned prefix residues are 1,3,1,3. At the
third bit, the provisional six is reduced to one; the final bit then gives
three. The residue invariant explains why each append needs at most one
subtraction, while the full implementation must also retain the extra carry.
Resume to reach each new visit to the same source instruction.

In \texttt{worked/modular-product.asm}, multiplier \texttt{11} selects two
contributions from base 250 modulo 251. The first residue is 250; the doubled
base residue is 249; their sum reduces to 248. Before using the modular adder
inside your multiplier, establish its reduced-input precondition. The multiplier
itself accepts arbitrary inputs. The multiplier's caller has not promised reduced operands.

\subsection*{Shift the whole exponent and preserve the outer calculation}

\texttt{rsa\_pow128(a0 base,a1 exponent,a2 N,a3 output)} uses the same power
invariant as the scalar routine, with modular multiplication replacing the
single-register product. All 128 exponent bits participate. Scanning only its
first word or first doubleword silently changes high exponents to smaller ones.

A whole-number right shift must bring bits across storage boundaries. If two
registers hold the low and high 64-bit halves, one possible sequence is:

\MARSCodeLabel{Move the high half's low bit into the low half}
\begin{minted}{mars-idea}
    dsrl $t0, $t0, 1
    dsll $t2, $t1, 63
    or $t0, $t0, $t2
    dsrl $t1, $t1, 1
\end{minted}

For exponent $2^{64}+1$, the initial halves are one and one. After shifting,
the low half must be \texttt{8000000000000000} and the high half zero.
Shifting each half independently would instead produce a zero exponent.
For base three modulo 19, the complete exponent $2^{64}+1$ gives 13, and
$2^{127}$ gives nine. A truncated exponent can accidentally pass a test whose
true result is one, so use a case that distinguishes the two calculations.

Keep the outer result, current base, remaining exponent and modulus address
across every nested helper call. Arguments and temporary registers are volatile.
A saved return address preserves where execution resumes, not the caller's
unfinished arithmetic state. Save and restore full preserved registers, including
their upper halves, and keep local temporary arrays separate from saved addresses.

The helpers allow exact sharing among read-only inputs, so squaring may pass the
same input address twice. Modular-product and power outputs remain disjoint from
their inputs. Only addmod has the additional output-alias permission. Establish
those relationships before reusing a temporary array for another role.

The native \texttt{worked/exponent-halves.asm} performs one logical shift
of a supplied 128-bit exponent. Its old high half is
\texttt{8000000000000001} and its low half is one. After shifting, the new low
half is \texttt{8000000000000000}; a bit has crossed the doubleword boundary.
The new high half is \texttt{4000000000000000}. Inspect that transfer before
choosing a termination test for the complete exponent scan. The experiment also
reduces one modulo one to zero, making the zero-exponent identity explicit.

\subsection*{Check the modulus before following other pointers}

For every modular multiword entry, first read its four-limb modulus and determine
whether all limbs are zero. On zero, return status one before reading other caller
inputs or accessing output. Those other pointers can be deliberately invalid in
a test. Checking only the low modulus limb would incorrectly reject a valid
modulus such as $2^{32}$.

For positive modulus one, the result is zero, including an exponent-zero power.
For a positive larger modulus and exponent zero, the result is one. These are
successful arithmetic results, distinct from the zero-modulus error.

The scalar operands of \texttt{rsa\_pow32} are values, not addresses; its output
pointer still has the no-access promise on zero modulus. The gcd entry has no
modulus error path. Write each function's arguments and return convention beside
your test cases rather than applying one function's rules to all seven.

\begin{figure}[!ht]
  \centering
  \begin{tikzpicture}[alt={Actual native Memory row after zero-modulus rejection. All sixteen selected output bytes remain a5 and are outlined. The separately checked status register is outside this crop.}]
    \node[anchor=north west,inner sep=0,font=\small\bfseries] (title) at (0,0) {Zero modulus leaves the output unchanged};
    \node[anchor=north west,inner sep=0] (a) at ([yshift=-2mm]title.south west) {\MARSRevisionCrop{revision/captures/rsa/rsa-invalid.png}{0.99\linewidth}{0.03125}{0.33888889}{0.58203125}{0.60888889}};
    \begin{scope}[shift={(a.south west)},x={($(a.south east)-(a.south west)$)},y={($(a.north west)-(a.south west)$)}]
      \draw[black,line width=.7pt] (.20,.01) rectangle (.965,.467);
    \end{scope}
  \end{tikzpicture}
  \caption{All sixteen selected modular-output bytes remain a5 after rejecting modulus zero. The outline encloses the selected interval. Separate checks verify status one and no caller writes; the status register is not pictured.}
  \label{fig:revision-rsa-native-invalid}
\end{figure}


Success writes exactly four bytes for pow32, sixteen for modular multiword
results and 32 for the full product. Failure preserves all caller objects.
An unchanged final value alone does not prove that an illegal write never occurred:
a store of the same byte is still a write. The public checks observe access
intervals as well as final bytes.

In \texttt{worked/validation.asm}, the modulus is readable but zero, and
the other pointer arguments contain the invalid address one. The helper reads
all four modulus limbs and returns status one. Main then finds the separate
output word unchanged. The independent checks observe caller accesses as well
as bytes: an unchanged result alone would not prove that no forbidden access
occurred. Keep caller publication outside the completed callee's ownership
boundary.

\subsection*{Use a small RSA example as an arithmetic test}

The last native assessment groups exercise the actual caller and different
operations within one simulator. The repeated-power case feeds one result into
the inverse calculation and then repeats the original request. State left over
from a previous call must not replace current arguments. After completing all
entries, run the supplied mode-zero caller to termination, then use the final
assessment to recheck every interface on the completed source. Later edits can
invalidate an earlier pass. Choose another input and explain its expected result
from the arithmetic invariant before testing it.

For optional standalone Debug observations, keep \texttt{mode} at its supplied
default zero in the source passed to the preparation tool. To observe the scalar power in a fresh source from the book
checkout, use:

\Needspace{6\baselineskip}
\begin{Verbatim}[fontsize=\small]
python3 packages/027-rsa-arithmetic/prepare.py s06-1 \
  --source revision/labs/files/rsa-arithmetic.asm \
  --output /private/tmp/rsa-word-observation.asm
\end{Verbatim}

\Needspace{9\baselineskip}
Replace the source path with your completed working source when testing your
implementation. The tool installs data and mode in the output copy, refuses an
existing output path and retains your instruction text. Recipes
\texttt{s11-1}, \texttt{s18-1}, \texttt{s29-1} and \texttt{s36-1} select the
full product, carry-sensitive addition, modular product and full-width power
examples. An unfinished entry still stops with \texttt{break}; prepared inputs
do not supply its implementation.

For a separate scalar reverse run, set the first words of \texttt{left} and
\texttt{right} to 2790 and 2753, keep modulus 3233 and select mode one.
The full-width power caller uses mode six and reads its four-limb exponent from
\texttt{exponent} instead of \texttt{right}. The function's argument registers,
not the fixture labels themselves, determine which values it receives.

For standalone CLI checking, use the nine original public scripts listed in
\texttt{packages/027-rsa-arithmetic/package.json}. Their names separate gcd,
word power, product, addition, reduction, modular product and full-width power,
plus caller and repeated-call behavior. Use
\texttt{java -jar /path/to/cli.jar grade CHECK SOURCE}. Keep the supplied caller
when using its checks; private function labels and exact instruction counts are
not part of free-writing acceptance.

Build an independent test set that includes zero and one moduli, zero exponents,
top-bit-set scalar products, carries and borrows across every limb, maximum
full products, high exponent bits, exact addmod aliases and moved input/output
slices. Compare values with an arbitrary-precision calculation, then check
preservation and ownership separately. A round trip alone can hide two matching
mistakes.

The example explains RSA's modular-power arithmetic. A complete encryption or
signature scheme additionally defines how messages are encoded and checked.
RFC 8017 describes OAEP for an encryption scheme and PSS for a signature scheme;
this exercise implements neither encoding. Its four-limb numbers are chosen so
we can inspect the arithmetic, not as a usable RSA key size. Nor do the data-dependent
loops establish constant-time execution. Keep those scheme and implementation
questions separate from the numerical and memory contracts you have tested here.
